\documentclass[presentation,10pt]{beamer}
\usetheme{Madrid}
\usecolortheme{whale}

% --- Packages ---
\usepackage{amsmath}
\usepackage{amsfonts}
\usepackage{amssymb}
\usepackage{booktabs}
\usepackage{graphicx}

% --- Presentation Metadata ---
\title[Subsonic Marshak Wave]{Complete Self-Similar Solution of the Subsonic Marshak Wave}
\subtitle{Theory and Numerical Validation Structure}
\author[Shussman \& Heizler]{Tomer Shussman\inst{1} \and Shay I. Heizler\inst{2}}
\institute[TAU / BIU]{
  \inst{1}Tel Aviv University / Soreq Nuclear Research Centre \and % [cite: 2]
  \inst{2}Bar-Ilan University / Nuclear Research Centre Negev % [cite: 2]
}
\date{Physics Conference 2026}

\begin{document}

% --- Title Slide ---
\begin{frame}
  \titlepage
\end{frame}

% --- Slide 1: Motivation ---
\begin{frame}{Motivation}
  \begin{block}{Inertial Confinement Fusion (ICF)}
    \begin{itemize}
      \item **Hohlraum Modeling:** Crucial for radiation temperature prediction and experimental measurement accuracy[cite: 5, 6].
      \item **Energy Delivery:** Key to analyzing laser--walls coupling and refining pulse shaping techniques[cite: 7, 8].
      \item **Historical Context:** Builds upon classical foundations established by \emph{Lindl (1995)}, \emph{Kauffman (1995)}, and \emph{Rosen (2005)}[cite: 9].
    \end{itemize}
  \end{block}

  \begin{block}{Equation of State (EOS) Verification}
    \begin{itemize}
      \item Critical for laboratory astrophysics and high-energy density physics (HEDP) experiments under dynamic, time-varying pressure profiles \emph{[Thomas 2006]}[cite: 10].
    \end{itemize}
  \end{block}
\end{frame}

% --- Slide 2: Physical Mechanism ---
\begin{frame}{The Subsonic Wave Mechanism}
  \begin{itemize}
    \item We treat a semi-infinite wall subjected to high-temperature boundary heating[cite: 43].
    \item Dynamic equilibrium is mediated strictly via **photon transport**[cite: 44].
    \item In the optically thick limit, an exceptionally sharp temperature front develops[cite: 45].
  \end{itemize}

  \vspace{0.3cm}
  \begin{alertblock}{Subsonic Regime Conditioning}
    If the localized heat front travels **slower than the material speed of sound**, hydrodynamics decouples into two distinct zones: a **shock region** driven forward ahead of the high-temperature **ablation region**[cite: 46].
  \end{alertblock}
\end{frame}

% --- Slide 3: Methodological Approach ---
\begin{frame}{Our Work -- Overview \& Strategy}
  \begin{itemize}
    \item **Decoupled Patching:** Instead of a single complex hydro-transport fit, we solve the ablation zone and the shock zone independently, later matching them via self-consistent boundary logic[cite: 69, 471].
    \item **Self-Similarity:** In each localized frame, we systematically neglect one governing dimensional parameter to yield exact self-similar ODEs[cite: 70].
    \item **Inter-Zone Coupling:** The thermodynamic properties calculated at the heat front serve directly as the driving boundary condition for the downstream shock equations[cite: 71].
  \end{itemize}

  \begin{table}
    \centering
    \begin{tabular}{lll}
      \toprule
      \textbf{Parameter Pool} & \textbf{Ablation Region} & \textbf{Shock Region} \\
      \midrule
      Density Profile & Assumed $\rho \to 0$ at source [cite: 95, 98] & Captures background $\rho_0$ [cite: 142] \\
      Radiation Flux ($B$) & Fully Coupled [cite: 90, 94] & Neglected ($B \approx 0$) [cite: 141] \\
      \bottomrule
    \end{tabular}
    \caption{Strategic scaling assumptions across regions.}
  \end{table}
\end{frame}

% --- Slide 4: Governing Equations ---
\begin{frame}{Statement of the Problem}
  The 1D radiation-hydrodynamic transport framework in Lagrangian mass coordinates $m$ is stated as[cite: 86]:
  
  $$\frac{\partial V}{\partial t} = \frac{\partial u}{\partial m}$$ % [cite: 87]
  $$\frac{\partial u}{\partial t} = -\frac{\partial P}{\partial m}$$ % [cite: 88]
  $$\frac{\partial e}{\partial t} + P\frac{\partial V}{\partial t} = \frac{\partial}{\partial m}\left(\frac{c}{3\kappa_R}\frac{\partial(aT^4)}{\partial m}\right)$$ % [cite: 88]

  \begin{block}{Combined Nonlinear Energy Equation}
    By tracking power-law material properties, the energy diffusion scales as[cite: 90]:
    $$\frac{1}{r}\frac{\partial PV}{\partial t} + P\frac{\partial V}{\partial t} = B\frac{\partial}{\partial m}\left(V^\lambda \frac{\partial}{\partial m}(PV^{1-\mu})^{\frac{4+\alpha}{\beta}}\right)$$ % [cite: 90]
    where the structural constant is $B = \frac{16\sigma}{3(4+\alpha)}g(rf)^{\frac{-(4+\alpha)}{\beta}}$[cite: 90].
  \end{block}
\end{frame}

% --- Slide 5: Material & Boundary Assumptions ---
\begin{frame}{Thermodynamic \& Material Modeling}
  To construct a closed analytic framework, variables are mapped using power-law dependencies[cite: 73]:
  
  \begin{columns}[t]
    \begin{column}{0.5\textwidth}
      \begin{block}{Boundary Heating Profile}
        $$T(t) = T_0 t^\tau$$ % [cite: 79]
      \end{block}
      \begin{block}{Rosseland Mean Opacity}
        $$\frac{1}{\kappa_R} = g T^\alpha \rho^{-\lambda}$$ % [cite: 80]
      \end{block}
    \end{column}
    
    \begin{column}{0.5\textwidth}
      \begin{block}{Specific Internal Energy}
        $$e = f T^\beta \rho^{-\mu}$$ % [cite: 81]
      \end{block}
      \begin{block}{Equation of State (EOS)}
        $$P = r\rho e \equiv (\gamma - 1)\rho e$$ % [cite: 82]
      \end{block}
    \end{column}
  \end{columns}
  
  \vspace{0.4cm}
  Background density is initialized uniformly to $\rho_0$ ahead of the wavefront[cite: 78].
\end{frame}

% --- Slide 6: Ablation Self-Similar Scaling ---
\begin{frame}{The Ablation Region Properties}
  The ablation zone is completely characterized by four dimensional parameters: $B, T_0, \rho_0 (V_0),$ and $t$[cite: 94]. Transforming to the dimensionless similarity coordinate $\xi$ scales the structural variables[cite: 101]:

  \begin{itemize}
    \item **Lagrangian Coordinate Scaling:** [cite: 101]
    $$m = \xi \cdot \left(B^{2-\mu}T_0^{8+2\alpha-2\beta+\beta\lambda-(4+\alpha)\mu} t^{2+2(4+\alpha-\beta)\tau - \mu(3+(4+\alpha)\tau) + \lambda(2+\beta\tau)}\right)^{\frac{1}{4+2\lambda-4\mu}}$$
    \item **Local Structural Functions:** Evaluated using dimensional scaling tracks [cite: 102]
    $$X(m,t) = \tilde{X}(\xi) B^{w_{X1}} T_0^{w_{X2}} t^{w_{X3}}$$
  \end{itemize}

  \begin{block}{Boundary Conditions for Ordinary Differential Equations}
    Boundary bounds are enforced exactly at the zero-volume ablation edge $\xi_F$[cite: 106]:
    $$\tilde{V}(\xi_F)=0, \quad \left.\frac{\partial\tilde{V}}{\partial\xi}\right|_{\xi_F}=0, \quad \tilde{P}(0)=0, \quad \tilde{u}(\xi_F)=0$$ % [cite: 106, 107, 108]
  \end{block}
\end{frame}

% --- Slide 7: Shock Boundary Mechanics ---
\begin{frame}{The Shock Region Integration}
  Once the ablation scale maps the pressure at the backend boundary, the shock zone is localized via $P(t) = P_0 t^{\tau_S}$[cite: 139]. Neglecting active radiation transport ($B \to 0$) yields classic strong shock relations[cite: 140, 141]:

  \begin{block}{Rankine-Hugoniot Mass, Momentum, and Energy Conditions}
    $$\frac{D - u_S}{V_S} = \frac{D}{V_0}, \quad P_S = \frac{D u_S}{V_0}, \quad P_S u_S V_0 = D\left(e_S + \frac{u_S^2}{2}\right)$$ % [cite: 142, 143]
  \end{block}

  \begin{itemize}
    \item **EOS Divergence:** To retain generality, the material model inside the shock domain is permitted to utilize alternative polytropic definitions ($r \neq r'$)[cite: 144, 146].
    \item **Shock Coordinate Mapping:** mass variables scale via[cite: 180]:
    $$m_S = \xi P_0^{1/2} V_0^{-1/2} t^{1 + \frac{\tau_S}{2}}$$
  \end{itemize}
\end{frame}

% --- Slide 8: Simulation Setup Case 1 ---
\begin{frame}{Simulation Setup: Case 1 ($\tau = 0$)}
  \begin{block}{Medium Framework: Gold Target}
    The numerical validation uses an ideal gas with polytropic index $\gamma = \frac{5}{4}$ and initial high material density $\rho_0 = 19.32\text{ g/cm}^3$[cite: 477, 480].
  \end{block}

  \begin{columns}[t]
    \begin{column}{0.5\textwidth}
      \textbf{Specific Internal Energy:} [cite: 481]
      $$e(T,\rho) = 3.4\times 10^{13} \left(\frac{T}{\text{HeV}}\right)^{1.6} \left(\frac{\rho}{\text{g/cm}^3}\right)^{-0.14}\frac{\text{erg}}{\text{g}}$$
    \end{column}
    \begin{column}{0.5\textwidth}
      \textbf{Rosseland Opacity Coefficient:} [cite: 481]
      $$\kappa_R(T,\rho) = 7200 \left(\frac{T}{\text{HeV}}\right)^{-1.5} \left(\frac{\rho}{\text{g/cm}^3}\right)^{0.2}\frac{\text{cm}^2}{\text{g}}$$
    \end{column}
  \end{columns}

  \vspace{0.3cm}
  \begin{block}{Boundary Conditions \& Self-Similar Evolution ($\tau = 0$)}
    Surface drive is maintained via a fixed thermal profile $T_s(t) = 1\text{ HeV}$[cite: 479, 485]. 
    The dynamic radiation bath drive and the calculated mass front position trace as[cite: 488, 489]:
    $$T_{\text{bath}}(t) = \left[1 + 0.169141 \left(\frac{t}{\text{ns}}\right)^{-\frac{79}{192}}\right]^{1/4} T_s(t)$$ % [cite: 488]
    $$m_{\text{front}}(t) = 1.01889\times 10^{-3} \left(\frac{t}{\text{ns}}\right)^{\frac{33}{64}}\text{ g/cm}^2$$ % [cite: 489]
  \end{block}
\end{frame}

% --- Slide 9: Simulation Setup Case 2 ---
\begin{frame}{Simulation Setup: Case 2 ($\tau \approx 0.123$)}
  \begin{block}{Time-Independent Surface Flux Environment}
    To simulate a completely steady, time-independent surface radiation energy flux ($c_S = 0$), the temporal power component is optimized analytically[cite: 490, 494]:
    $$\tau = \frac{2-3\mu}{8+2\alpha+2\beta+3\beta\lambda-3(4+\alpha)\mu} = \frac{158}{1285} \approx 0.12296$$ % [cite: 490, 496]
  \end{block}

  \begin{itemize}
    \item **Material Invariance:** Identical baseline properties as Case 1 ($\gamma=\frac{5}{4}$, same $e$, same $\kappa_R$)[cite: 493].
    \item **Time-Dependent Surface Drive:** Scales upward explicitly as $T_s(t) = \left(\frac{t}{\text{ns}}\right)^{\frac{158}{1285}}\text{ HeV}$[cite: 496].
  \end{itemize}

  \begin{block}{Self-Similar Validation Metrics}
    Analytical tracking metrics for the flux-matched drive are implemented as[cite: 496, 497]:
    $$T_{\text{bath}}(t) = \left[1 + 0.204154 \left(\frac{t}{\text{ns}}\right)^{-\frac{632}{1285}}\right]^{1/4} T_s(t)$$ % [cite: 496]
    $$m_{\text{front}}(t) = 8.34319\times 10^{-4} \left(\frac{t}{\text{ns}}\right)^{\frac{103}{257}}\text{ g/cm}^2$$ % [cite: 497]
  \end{block}
\end{frame}

% --- Slide 10: Summary ---
\begin{frame}{Summary}
  \begin{itemize}
    \item **Complete Solution Matrix:** Provides a rigorous, closed-form analytic description of a subsonic radiation heat wave for generalized power-law boundaries[cite: 470].
    \item **Robust Decoupled Framework:** Separates the localized ablation and shock regions successfully, utilizing cross-boundary patching values[cite: 471].
    \item **Practical Utility:** Enables higher fidelity planning for laser-plasma and high-energy-density experiments while revealing the physical structure of ablation-driven shock waves[cite: 472].
  \end{itemize}

  \vfill
  \begin{center}
    \begin{alertblock}{Primary Literature Source}
      Full mathematical derivations and comprehensive simulation arrays are published in: \\
      \vspace{0.1cm}
      \centering{\textbf{Phys. Plasmas 22, 082109 (2015)}} [cite: 473]
    \end{alertblock}
  \end{center}
\end{frame}

\end{document}
